% !TeX root=../main.tex

\chapter{بحث و نتیجه‌گیری}
\label{chap:conclusion}

\section{مرور مسیر پروژه}

این پروژه از یک مسئلهٔ عملی در زنجیرهٔ تحلیل ربکا آغاز شد: فضای حالت زمان‌دار تولیدشده توسط افرا می‌تواند جزئیاتی داشته باشد که برای سطح مشاهدهٔ انتخابی کاربر ضروری نیستند، اما تحلیل‌گر همچنان باید همهٔ آن‌ها را نگهداری و پیمایش کند. تعریف شبیه‌سازی دوسویهٔ ضعیف زمان‌دار مبنایی برای نادیده‌گرفتن کنش‌های داخلی با حفظ رفتار قابل مشاهده و زمان فراهم می‌کند \cite{ReynierSangnier2009,LiebEtAl2025}. فاصلهٔ میان این تعریف و یک ابزار قابل استفاده شامل چند مسئلهٔ مهندسی و معنایی بود: قرارداد واقعی فایل افرا، هویت تعامل‌ها، تعبیر تأخیرهای چندمرحله‌ای، ساخت مدل خارج‌قسمتی و وارسی نتیجهٔ تحویلی.

کار اصلی پروژه پیاده‌سازی این فرایند در ابزار \lr{awtr} بود. ابزار فایل \lr{.statespace} را می‌خواند، تعامل‌های خارج از مجموعهٔ مشاهده را پنهان می‌کند، بخش دسترس‌پذیر مدل را نگه می‌دارد، تأخیرها را مطابق معنای انتخابی پالایش می‌کند و پس از ساخت رابطهٔ گذار ضعیف، افراز پایدار و مدل خارج‌قسمتی را به دست می‌آورد. وارسی‌گری جداگانه رابطهٔ مدل پالایش‌شده و خروجی را در هر دو جهت بررسی می‌کند. نتیجه همراه با نگاشت افراز، منشأ یال‌ها و سنجه‌های اجرا ذخیره می‌شود تا صرفاً به یک پاسخ مثبت بدون شاهد اکتفا نشود.

\section{پاسخ به پرسش‌های پروژه}

\subsection{خواندن و مستقل‌کردن خروجی افرا}

پرسش نخست دربارهٔ امکان مصرف مستقیم خروجی زمان‌دار افرا بود. بررسی تولیدکننده و خواننده‌های رسمی نشان داد فایل واقعی \lr{TTS} قراردادی دارد که تنها از روی طرح داده روشن نمی‌شود؛ از جمله، عنصر ریشه ممکن است در انتهای فایل بسته نشده باشد و نخستین حالت نقش حالت آغازین را داشته باشد. خوانندهٔ پیاده‌سازی‌شده این موارد را بدون تغییر دستی فایل مدیریت و داده را پشت رابطی مستقل از افرا به مدل دامنه تبدیل می‌کند. آزمون قرارداد نشان می‌دهد منبع فایل و منبع درون‌حافظه‌ای به نتیجهٔ یکسان می‌رسند. بنابراین، برای دامنهٔ \lr{TTS} و فیلدهای مستندشده، پاسخ پرسش نخست مثبت است. این نتیجه به \lr{FTTS} تعمیم داده نمی‌شود.

\subsection{کاهش همراه با وارسی مستقل}

پرسش دوم ساخت یک مدل کوچک‌تر و وارسی رابطهٔ آن با ورودی بود. هر چهار ورودی سنجه‌ای و هر هفت فضای حالت تولیدشده از مدل‌های ربکا به خروجی کوچک‌تری رسیدند و وارسی‌گر مستقل همهٔ رابطه‌های تحویلی را پذیرفت. برای سنجش اینکه پذیرش صرفاً نتیجهٔ یک وارسی‌گر همیشه‌مثبت نیست، پنج نوع تغییر مخرب روی مدل خارج‌قسمتی اعمال شد و وارسی‌گر همه را رد کرد. ۱۱۶ آزمون \lr{JUnit} و ۶۲ بررسی انتهابه‌انتها نیز بدون شکست اجرا شدند. این شواهد از درستی خروجی‌های موجود پشتیبانی می‌کنند، اما جای اثبات درستی خود پیاده‌سازی یا بازبینی مستقل تعریف را نمی‌گیرند.

واژهٔ «کاهش» در این نتیجه باید دقیق تفسیر شود. مدل خروجی رفتاری را که مجموعهٔ مشاهده و معنای زمان تعیین می‌کنند حفظ کرده است؛ با تغییر مجموعهٔ مشاهده، افراز نیز می‌تواند تغییر کند. همچنین، ابزار کوچک‌ترین مدل ممکن یا یک نمایش کانونی را تضمین نمی‌کند. ادعای پشتیبانی‌شده این است که خارج‌قسمت‌های تولیدشده برای نمونه‌های ارزیابی از ورودی کوچک‌تر و مطابق وارسی‌گر، با آن دارای شبیه‌سازی دوسویهٔ ضعیف زمان‌دار بوده‌اند.

\subsection{اثر معنای زمان و اندازهٔ کاهش}

پرسش سوم به اثر شکستن یال‌های تأخیر و مقدار کاهش مربوط بود. در حالت پیش‌فرض \lr{unit}، هر تأخیر چندواحدی به گام‌های یک‌واحدی شکسته می‌شود؛ بنابراین یک تأخیر ده‌واحدی می‌تواند با اجرای $7+\tau+3$ یا $3+\tau+7$ تطبیق یابد. سه بیان مطالعهٔ موردی خانهٔ هوشمند در این حالت هم‌ارز شناخته شدند و پس از فشرده‌سازی به ساختاری با ۱۰ حالت و ۱۲ گذار رسیدند؛ جفت \lr{mood1} و \lr{mood2} نیز، که از دو مدل واقعی ربکا تولید شده‌اند، هم‌ارز شناخته شدند. در حالت \lr{strict}، که یال‌ها را دست‌نخورده نگه می‌دارد، نه نمونهٔ پایهٔ خانهٔ هوشمند با \lr{tc2step} هم‌ارز است و نه \lr{mood1} با \lr{mood2}. این تفاوت، چنان‌که بخش \ref{sec:unit-splitting} توضیح داد، نه یک انتخاب سلیقه‌ای بلکه نتیجهٔ کنارگذاشتن ویژگی پیوستگی است: بدون شکستن یال‌ها، الگوریتم روی ساختاری کار می‌کند که ویژگی‌های استاندارد یک سامانهٔ انتقال زمان‌دار را ندارد و آنچه محاسبه می‌کند دیگر شبیه‌سازی دوسویهٔ ضعیف زمان‌دار به معنای تعریف \ref{def:weak-timed-bisimulation} نیست.

در ورودی کوچک، تعداد حالت‌ها از ۶ به ۲ و تعداد گذارها از ۷ به ۲ کاهش یافت. برای سه نمونهٔ خانهٔ هوشمند، نسبت کاهش حالت به‌ترتیب \lr{37.5}، \lr{60.0} و \lr{76.2} درصد و نسبت کاهش گذار \lr{33.3}، \lr{57.1} و \lr{76.9} درصد بود؛ همهٔ این ارقام با مجموعهٔ مشاهدهٔ \lr{getSense}، \lr{activateh} و \lr{switchoff} معنا دارند. در فضای حالت‌های تولیدشده از ربکا، نسبت کاهش حالت میان \lr{16.7} و \lr{89.5} درصد بود و در هر سه دسته، مدل‌های هم‌ارز به مدل خارج‌قسمتی یکسانی رسیدند. افزایش نسبت کاهش در نمونه‌هایی که جزئیات داخلی بیشتری دارند قابل انتظار است، اما این تعداد نقطهٔ ارزیابی برای نتیجه‌گیری دربارهٔ روند مقیاس‌پذیری کافی نیست.

\section{تفسیر تصمیم‌های طراحی}

پالایش تأخیرها به گام‌های واحد، تعریف جمع‌پذیر زمان را به مسئله‌ای متناهی با یک برچسب زمانی واحد تبدیل می‌کند. مزیت آن سادگی و امکان استفاده از همان پالایش افراز برای کنش و زمان است. هزینهٔ آن افزودن $d-1$ حالت برای یک یال با مدت $d$ است. فشرده‌سازی زنجیره‌های صرفاً زمانی پس از ساخت خارج‌قسمت، اندازهٔ مدل تحویلی را کاهش می‌دهد، ولی هزینهٔ نمایش میانی را حذف نمی‌کند. بنابراین این تصمیم برای ورودی‌های فعلی مناسب بوده است، نه لزوماً برای مدل‌هایی با ثابت‌های زمانی بسیار بزرگ.

جداسازی وارسی‌گر از کاهش‌دهنده نیز یک مصالحهٔ مهندسی است. وارسی‌گر بستارها و تعهدهای رابطه را با ساختار داده و مسیر کد متفاوتی محاسبه می‌کند؛ ازاین‌رو، بخشی از خطاهای شاخص‌گذاری، ساخت امضا یا فشرده‌سازی می‌تواند آشکار شود. بااین‌حال، هر دو مؤلفه تعبیر معنایی مشترکی دارند و خطای مشترک در همان تعبیر ممکن است از هر دو عبور کند. آزمون جهش اعتماد به توان ردکردن خروجی معیوب را افزایش می‌دهد، اما اثبات صوری وارسی‌گر نیست.

خروجی‌های چندگانه نیز دو نیاز متفاوت را پاسخ می‌دهند. \lr{reduced.statespace} برای مشاهده در ابزارهای افرا سودمند است، ولی قالب آن همهٔ منشأها و عضویت رده‌ها را نگه نمی‌دارد. در مقابل، \lr{reduced.json} نتیجهٔ اصلی و بدون اتلاف پروژه است و \lr{partition.json} امکان ردیابی حالت ورودی تا حالت یا یال خروجی را فراهم می‌کند. این تفکیک اجازه می‌دهد سازگاری با ابزار موجود، اطلاعات لازم برای بازتولید نتیجه را محدود نکند.

\section{جایگاه دستاورد}

این پایان‌نامه تعریف تازه‌ای از شبیه‌سازی دوسویه یا الگوریتم تازه‌ای برای محاسبهٔ آن ارائه نمی‌کند. تعریف رفتاری از برار و همکاران و ساختار الگوریتم از پالایش افراز گروته و واندراخر گرفته شده است \cite{BerardEtAl2008,GrooteVaandrager1990}. سهم پروژه، تبدیل این مبنا به یک مسیر اجرایی و آزموده برای مدل معنایی ربکاست: قرارداد ورودی افرا از کدهای رسمی استخراج شده، نحوهٔ برخورد با تأخیرها از ویژگی پیوستگی سامانهٔ انتقال زمان‌دار استخراج و به‌صورت گزینه‌ای صریح پیاده‌سازی شده، مدل خارج‌قسمتی همراه با منشأ ساخته شده و نتیجه با پیاده‌سازی جداگانه و آزمون‌های منفی بررسی شده است.

پیش از این، کاربر زنجیرهٔ ربکا و افرا ابزاری در اختیار نداشت که فضای حالت زمان‌دار تولیدشده را نسبت به مجموعه‌ای از تعامل‌های قابل مشاهده کاهش دهد؛ وجود چنین ابزاری، مستقل از اندازهٔ کاهشی که روی نمونه‌های خاص به دست می‌آید، خودِ دستاورد این کار است \cite{AfraWebsite2026,RMCGitHub2026,AfraWeakTimedReduction2026}.

از دید کاربردی، ابزار موجود یک نمونهٔ مستقل از خط فرمان است و هنوز جزئی از محیط افرا نیست. در نتیجه، دستاورد فعلی را باید زیرساختی برای ارزیابی روش و ادغام آینده دانست، نه بهبود اثبات‌شدهٔ کارایی خود افرا. داده‌های فصل \ref{chap:results} نشان می‌دهند کاهش روی نمونه‌های موجود رخ داده و هزینهٔ اجرا در محیط ثبت‌شده کم بوده است؛ اما چون زمان تحلیل بعدی با افرا اندازه‌گیری نشده، نمی‌توان از این نتایج دربارهٔ شتاب نهایی وارسی مدل ادعای کمی مطرح کرد.

\section{محدودیت‌ها}

محدودیت‌های اصلی نتیجه عبارت‌اند از:

\begin{itemize}
\item بزرگ‌ترین ورودی پیش از پالایش زمانی ۴۲ حالت دارد و نمایندهٔ یک مدل صنعتی نیست؛
\item ارزیابی چهار ورودی سنجه‌ای و هفت فضای حالت تولیدشده از مدل‌های ربکای نوشته‌شده برای همین ارزیابی را پوشش می‌دهد و هر زمان اجرا یک‌بار ثبت شده است؛
\item فایل افرا از روی نمایش تولیدشده بازسازی شده و خروجی مستقیم یک اجرای رسمی در دسترس نبوده است؛
\item پشتیبانی زمانی به \lr{TTS}، زمان گسسته و تأخیرهای صحیح محدود است؛
\item پالایش واحدی می‌تواند برای تأخیرهای بزرگ نمایش میانی بسیار بزرگی بسازد؛
\item نتیجه‌های هم‌ارزی به قرارداد معنایی تأییدشدهٔ جمع‌پذیری زمان وابسته‌اند و به خوانش \lr{strict} تعمیم داده نمی‌شوند؛
\item خارج‌قسمت وارسی می‌شود، اما مینیمال یا کانونی‌بودن آن اثبات نشده است؛
\item سنجش حاضر کاهش فضای حالت را گزارش می‌کند و اثر آن بر زمان یک وارسی مدل کامل را اندازه نمی‌گیرد.
\end{itemize}

این محدودیت‌ها نتیجهٔ اصلی روی نمونه‌های ثبت‌شده را نقض نمی‌کنند، اما مرز تعمیم آن را تعیین می‌کنند. به‌ویژه، درصدهای کاهش تنها همراه با مجموعهٔ مشاهده و معنای زمانی همان اجرا قابل تفسیرند.

\section{کارهای آینده}

نخستین ادامهٔ طبیعی، افزودن آداپتور درون‌فرایندی و فرمان متناظر در افراست تا مدل بدون فایل واسط به هستهٔ کاهش تحویل شود. پیش از آن، نمونهٔ بازسازی‌شده باید با یک خروجی مستقیم و قابل انتشار افرا جایگزین و قرارداد خواننده روی چند مدل واقعی دیگر بررسی شود. تشخیص صریح \lr{FTTS} نیز باید افزوده شود تا ابزار به‌جای تفسیر خاموش یک قالب پشتیبانی‌نشده، خطای روشن ارائه کند.

در سطح معنایی، جمع‌پذیری زمان به‌عنوان قرارداد اصلی روش تثبیت شده است. مقایسهٔ تفاضلی با یک پیاده‌سازی مستقل دیگر یا صورت‌بندی مرجع می‌تواند خطر خطای مشترک میان کاهش‌دهنده و وارسی‌گر را کمتر کند. اثبات برخی ناورداهای کلیدی، مانند حفظ رابطه در فشرده‌سازی زنجیرهٔ زمان، گام بعدی دقیق‌تری از افزودن آزمون‌های مشابه است.

برای مقیاس بزرگ‌تر، می‌توان از پالایش افراز کارآمدتر—برای نمونه صورت $O(m\log n)$ که برای شبیه‌سازی دوسویهٔ انشعابی ارائه شده است \cite{GrooteEtAl2017}—و نمایش نمادین زمان با ناحیه یا زون استفاده کرد تا نیاز به ساخت یک حالت برای هر واحد تأخیر کاهش یابد. ارزیابی این نسخه باید مجموعه‌ای متنوع از مدل‌های ربکا، تأخیرهای بزرگ‌تر و تعداد اجرای تکراری داشته باشد و علاوه بر اندازهٔ خارج‌قسمت، زمان و حافظهٔ کل جریان وارسی مدل را پیش و پس از کاهش مقایسه کند.

مسیر دیگری نیز وجود دارد که ارزش بررسی دارد و هدف آن رسیدن به همان نتیجه بدون شکستن یال‌های تأخیر است. بخش \ref{sec:unit-splitting} نشان داد که اگر یال‌ها دست‌نخورده بمانند، برخی حالت‌ها—مانند حالت میانی $q_1$ در مثال آن بخش—در طرف مقابل شریکی پیدا نمی‌کنند و ساخت رابطه شکست می‌خورد. به‌جای ساختن آن شریک با شکستن یال، می‌توان تعریف رابطه را اندکی سست کرد و اجازه داد تعدادی حالت بی‌شریک باقی بماند: حالت‌هایی که صرفاً نقطه‌ای میانی روی یک انتظار پیوسته‌اند، هیچ کنش قابل مشاهده‌ای از آن‌ها آغاز نمی‌شود و هیچ انتخاب تازه‌ای نمی‌سازند. چنین حالتی را می‌توان بدون آسیب به رفتار قابل مشاهده از الزام تطبیق معاف کرد و الگوریتم را روی همان گراف ورودی، یعنی در همان حالت \lr{strict}، اجرا کرد.

کار لازم برای این مسیر، صورت‌بندی دقیق شرطی است که این حالت‌های بی‌شریک را مشخص می‌کند و سپس اثبات اینکه رابطهٔ حاصل هنوز یک شبیه‌سازی دوسویهٔ ضعیف زمان‌دار به معنای تعریف \ref{def:weak-timed-bisimulation} است. اگر این اثبات به دست آید، نتیجه الگوریتمی است که به همان هم‌ارزی می‌رسد اما هزینهٔ حافظهٔ پالایش واحدی—یعنی افزودن $d-1$ حالت برای هر یال با مدت $d$—را نمی‌پردازد. این مسیر مستقل از راه‌حل نمادین بند پیش است و می‌تواند در زمان گسسته نیز به کار رود.

در نهایت، قالب خروجی افرا ظرفیت نگهداری همهٔ منشأهای یک یال خارج‌قسمتی را ندارد. تعریف یک افزونهٔ سازگار برای این فراداده یا نمایش مستقیم آن در رابط افرا می‌تواند قابلیت ردیابی \lr{JSON} را به جریان تصویری منتقل کند و بررسی دستی علت ادغام حالت‌ها را آسان‌تر سازد.

\section{جمع‌بندی نهایی}

نتایج نشان می‌دهند تعریف شبیه‌سازی دوسویهٔ ضعیف زمان‌دار را می‌توان برای خروجی \lr{TTS} افرا به یک فرایند کاهش قابل اجرا و قابل ردیابی تبدیل کرد. ابزار ساخته‌شده در هر یازده ورودی ارزیابی—چهار ورودی سنجه‌ای و هفت فضای حالت تولیدشده از ربکا—مدلی کوچک‌تر تولید کرد، رابطهٔ خروجی‌ها در وارسی مستقل پذیرفته شد و آزمون‌های منفی توانستند تغییرهای معیوب را آشکار کنند. مطالعهٔ موردی همچنین نشان داد یک تصمیم ظاهراً کوچک دربارهٔ ترکیب تأخیرها می‌تواند پاسخ هم‌ارزی را تغییر دهد؛ بنابراین، ثبت معنای زمان به‌اندازهٔ خود الگوریتم اهمیت دارد.

این نتیجه یک گام مهندسی برای پیوند مبانی نظری و زنجیرهٔ ابزار ربکا و افراست. شواهد موجود برای درستی نمونه‌های بررسی‌شده و امکان کاهش آن‌ها کافی‌اند، اما برای ادعای مقیاس‌پذیری، مینیمال‌بودن خروجی یا بهبود کارایی افرا کافی نیستند. ادامهٔ کار باید بر تثبیت معنای زمان، ورودی‌های واقعی‌تر، ادغام با افرا و ارزیابی در مقیاس بزرگ‌تر متمرکز شود.
